\begin{table}[t]
\centering
\small
\caption{Q3 weather-intervention response fidelity on the frozen hot--dry
subset ($n=84$ pairs, 31 geographic clusters). Positive $\Delta$ loss means
that replacing actual future weather worsens the endpoint forecast.}
\label{tab:q3_weather_intervention}
\begin{tabular}{lcccc}
\toprule
Future-weather input & $R^2 \uparrow$ & RMSE $\downarrow$ &
$\Delta$ loss vs.\ actual & Geo-cluster 95\% CI \\
\midrule
Actual weather        & 0.6254 & 0.1492 & 0       & -- \\
Matched donor weather & 0.5893 & 0.1584 & 0.00257 & [0.00112, 0.00399] \\
Mean weather          & 0.5430 & 0.1971 & 0.01126 & [0.00547, 0.01708] \\
\bottomrule
\end{tabular}

\vspace{2pt}
\parbox{0.98\linewidth}{\footnotesize
Endpoint response-fidelity criterion: \textsc{Pass} for both matched-donor and
mean-weather replacements (10,000 bootstrap replicates). These scores are for
the 84-pair extreme subset, not the full temporal-shift test set. The separate
hot--dry enhancement interaction did \emph{not} pass:
$\Delta$ loss $=0.00044$, geo-cluster 95\% CI
$[-0.00216,\,0.00320]$. Thus the table supports weather-sensitive forecast
behavior, not an extreme-specific enhancement or causal counterfactual claim.}
\end{table}
